\documentclass[11pt,a4paper]{article}
\usepackage[utf8]{inputenc}\usepackage[T1]{fontenc}
\usepackage[margin=1in]{geometry}
\usepackage{amsmath,amssymb,booktabs,graphicx,enumitem,xcolor,parskip}
\usepackage[hidelinks]{hyperref}
\usepackage{titlesec}
\definecolor{qnavy}{RGB}{20,40,75}\definecolor{qgrey}{RGB}{90,90,90}\definecolor{qred}{RGB}{150,30,30}
\titleformat{\section}{\large\bfseries\color{qnavy}}{\thesection}{0.6em}{}
\titleformat{\subsection}{\bfseries\color{qnavy}}{\thesubsection}{0.5em}{}

\title{\bfseries\color{qnavy} Site identity is learnable from whole-slide foundation-model
features, and it inflates apparent performance by seventeen points in the encoder we started with}

\author{Reza Nehzati, Ph.D.\\[2pt]
{\small\itshape VMC MAR COM Inc.\ DBA HeyDonto, Knoxville, TN, United States}}
\date{15 August 2026}

\begin{document}
\maketitle
\thispagestyle{empty}

\begin{abstract}
\noindent
Whole-slide foundation models are increasingly proposed as the substrate for computational
companion diagnostics. We show that features from one such model (Phikon-v2) encode the
\emph{provenance} of a slide strongly enough to be recovered by a linear probe, and that this
provenance signal silently inflates cross-validated performance when training and evaluation
share contributing sites. Using 1{,}049 slides from 952 TCGA lung patients, we hold the data,
the architecture and every hyperparameter fixed and change only the cross-validation fold
assignment: grouping folds by tissue-source site gives subtype AUROC 0.799, while random folds
give 0.970 --- an inflation of \textbf{0.171} attributable to fold assignment alone.
\textcolor{qred}{Across five encoders this inflation ranges from 0.022 to 0.203 (relative
0.444--0.852), so 0.171 characterises \emph{this} encoder and the \emph{direction} is what
generalises; see \S Limitations.} The
random-fold figure reproduces an independently obtained centralized benchmark to within 0.0003,
establishing that the site-grouped number reflects honest out-of-site performance rather than a
weakened pipeline. Site identity is recoverable at balanced accuracy 0.706 across five
site-anchored partitions within TCGA (chance 0.200, permutation $p = 0.0099$), and
\emph{perfectly} between two genuinely independent external cohorts under Phikon-v2 (balanced
accuracy 1.000, chance 0.500, permutation $p = 0.005$, null centred at 0.506), holding unchanged
when restricted to H\&E slides only; a second encoder trained on natural images reaches $0.918 \pm 0.025$ over 25
cross-validation splits on the same slides, so the signature is encoder-independent while its
totality is not. Across five encoders in total the separation is present in every one
(0.918--1.000) and saturates in three, so the external probe is a presence test rather than a
ranking. Federated averaging across five
simulated silos recovers centralized subtype performance without moving data (0.976 vs 0.970)
while transmitting 301\,MiB per run, an 88-fold reduction against moving the features. On a
genuinely external institution set, subtype transfers zero-shot at AUROC 0.979 and requires no
labelled cases, whereas EGFR mutation status remains weak and power-limited. We argue that
site-grouped evaluation should be the default reporting standard for this class of model, and
that single-institution validation figures should be read as upper bounds.
\end{abstract}

\section{Introduction}

The commercial case for computational pathology rests on a model trained at one set of
institutions performing at another. That assumption is rarely tested in the way deployment
would test it. Cross-validation over a pooled multi-institution cohort permits patients from
the same contributing site to appear in both training and evaluation folds, and if the features
carry any site-specific signal, the model may exploit it.

This is not a hypothetical concern for foundation-model features. Slide preparation, staining,
scanner optics and compression differ measurably between institutions, and a representation
trained to be maximally informative about tissue will not have been trained to be invariant to
any of it.

We make three contributions. First, we isolate the magnitude of the inflation by changing
nothing but the fold assignment. Second, we show that the site signal is not an artefact of a
single collection: it is stronger between genuinely independent cohorts than within one.
Third, we quantify what federation buys and what it costs, and what a new institution must
supply to deploy.

\section{Data}

\subsection{Development cohort}

TCGA-LUAD and TCGA-LUSC diagnostic slides: 1{,}053 whole-slide images retrieved from the NCI
Genomic Data Commons, of which 1{,}049 encoded successfully, covering 952 patients across 67
tissue-source sites. Slides were tiled at approximately 0.5\,\textmu m/px with an Otsu
saturation tissue filter, yielding 14{,}721{,}903 tiles, and each tile encoded with Phikon-v2
(1024-dimensional).

\subsection{External institution set}

CPTAC-LUAD and CPTAC-LSCC: 1{,}343 slides from 432 patients, processed identically and held
back until the transfer experiments.

\subsection{Independent external cohorts}

A systematic sweep of all 3{,}411 open lung slide-image files in the Genomic Data Commons
identified only two cohorts independent of TCGA and CPTAC:

\begin{center}
\begin{tabular}{@{}llrl@{}}
\toprule
\textbf{Cohort} & \textbf{Setting} & \textbf{H\&E slides used} & \textbf{Note} \\
\midrule
CDDP\_EAGLE-1 & Italian population-based & 49 & all adenocarcinoma \\
CGCI-HTMCP-LC & HIV-positive, sub-Saharan & 29 & see below \\
\bottomrule
\end{tabular}
\end{center}

CGCI-HTMCP-LC ships 84 slides, but only 30 are haematoxylin and eosin; the remainder are
immunohistochemistry (p40, TTF-1, synaptophysin). Including them would allow a site probe to
separate this cohort on stain colour rather than provenance, so all non-H\&E slides were
excluded before any analysis. Four further slides, all immunohistochemistry, failed the tissue
filter and were excluded for the same reason. One H\&E slide failed retrieval. The usable
external set is therefore \textbf{78 H\&E slides from 78 patients}, and we treat it as a
qualitative test of whether the site signal generalises, not as a powered benchmark.

\section{Methods}

\subsection{Model}

A gated attention multiple-instance learning head over tile embeddings: a 1024$\rightarrow$512
projection, gated attention with a 256-dimensional bottleneck, and a linear output. Trained for
20 epochs with AdamW, learning rate $2\times10^{-4}$, weight decay $10^{-5}$, batch size 32,
tiles capped at 4{,}096 per slide. Patient-level predictions average over that patient's slides.

\subsection{The fold-assignment experiment}

The central comparison holds everything fixed except how folds are formed. Both arms use
identical patients (760 with matched molecular data, 413 adenocarcinoma and 347 squamous),
identical features, identical architecture, identical hyperparameters and identical seeds. The
arms differ only in the partition:

\begin{itemize}[leftmargin=1.3em,itemsep=2pt]
\item \textbf{Site-grouped}: \texttt{GroupKFold} on tissue-source site, five folds. No site
      contributes patients to both training and evaluation in any fold, asserted per fold.
\item \textbf{Random}: \texttt{StratifiedKFold}, five folds, sites free to appear on both sides.
\end{itemize}

We report the random arm not as a result but as a measurement of the inflation a conventional
analysis would have produced.

\subsection{Site probe}

A standardised logistic regression predicts the source of a slide from its mean tile embedding,
scored by balanced accuracy under stratified five-fold cross-validation, against a permutation
null of 200 label shuffles.

\subsection{Federation}

Five silos anchored on real tissue-source sites, compared under three regimes: each silo alone,
federated averaging with weights only crossing the boundary (10 rounds), and centralized
pooling. Five seeds per condition.

\section{Results}

\subsection{Fold assignment alone inflates subtype AUROC by 0.171}

\begin{center}
\begin{tabular}{@{}lcc@{}}
\toprule
& \textbf{Subtype AUROC} & \textbf{KEAP1 AUROC} \\
\midrule
Site-grouped (no shared sites) & 0.7992 & 0.6642 \\
Random folds (sites shared) & 0.9703 & 0.6660 \\
\midrule
\textbf{Inflation} & \textbf{+0.1710} & +0.0018 \\
\bottomrule
\end{tabular}
\end{center}

The random-fold subtype figure, 0.9703, reproduces an independently obtained centralized
benchmark of 0.970 to within 0.0003. This matters for interpretation: it establishes that the
site-grouped 0.799 is not the product of a weaker implementation but the honest out-of-site
performance of a pipeline whose capability is verified against the conventional number.

The asymmetry between endpoints is informative. Subtype is heavily inflated; KEAP1 mutation is
not inflated at all. \textcolor{qred}{(R18: this asymmetry reproduces under a second encoder ---
KEAP1 inflation is +0.0202 under \texttt{dinov2-large} against +0.0018 under Phikon-v2, still small
under a model chosen for an unrelated reason.)} Histological subtype prevalence varies between contributing institutions,
so site identity is predictive of subtype and a model can profit from recognising it. KEAP1
mutation status has no comparable institutional structure, and the shortcut is worth nothing.
Inflation is therefore not a property of the features alone but of the interaction between the
features and how the label is distributed across sites.

\subsection{Site identity is recoverable, and more so between real institutions}

\begin{center}
\begin{tabular}{@{}lrrrr@{}}
\toprule
\textbf{Probe} & \textbf{Bal.\ acc.} & \textbf{Chance} & \textbf{Null mean} & \textbf{Perm.\ $p$} \\
\midrule
Five site-anchored partitions within TCGA & 0.7062 & 0.2000 & 0.2029 & 0.0099 \\
EAGLE vs HTMCP, Phikon-v2 ($n=78$) & \textbf{1.0000} & 0.5000 & 0.5074 & 0.0050 \\
\quad \emph{--- H\&E only ($n=75$)} & \textbf{1.0000} & 0.5000 & 0.5058 & 0.0050 \\
\quad \emph{--- all H\&E in the archive ($n=76$)} & \textbf{1.0000} & 0.5000 & 0.5023 & 0.0050 \\
\quad \emph{--- \texttt{dinov2-large}, all H\&E ($n=76$)} & 0.8768 & 0.5000 & 0.4990 & 0.0050 \\
\quad \emph{--- \texttt{dinov2-large}, 25 CV splits} & $0.918{\pm}.025$ & 0.5000 & --- & --- \\
\quad \emph{--- \texttt{UNI}, 25 CV splits} & $0.994{\pm}.011$ & 0.5000 & --- & --- \\
\quad \emph{--- \texttt{H-optimus-0}, 25 CV splits} & $0.996{\pm}.008$ & 0.5000 & --- & --- \\
\quad \emph{--- \texttt{Virchow2}, 25 CV splits} & $0.960{\pm}.018$ & 0.5000 & --- & --- \\
TCGA vs EAGLE vs HTMCP (size-balanced) & 0.9816 & 0.3333 & 0.3345 & --- \\
\quad \emph{--- restricted to H\&E only} & 0.9814 & 0.3333 & 0.3660 & --- \\
\bottomrule
\end{tabular}
\end{center}

In every case the permutation null sits at chance, which is what distinguishes a real signal
from high-dimensional separability: with 1024 features and 75 external slides, a linear probe
could plausibly separate arbitrary labels, and the null establishes that it does not (null p95
0.617 against an observed 1.000).

\textcolor{qred}{\textbf{Corrected 2026-08-17 (R19).}} Three things in this table changed, and the
reason is worth stating rather than burying. The two external rows were originally computed in a
working session and transcribed here \emph{without a persisted artefact}, so they were not
verifiable; they now are. In re-running them we found that the input file --- named
\texttt{external\_he\_mean.npz} --- contained three HTMCP slides that are not H\&E (two chromogenic
ISH, one p16 immunohistochemistry), which for a probe scoring exactly 1.000 is the cheapest
available shortcut. Restricting to H\&E leaves the result \textbf{unchanged at 1.0000}, so the
H\&E-only rows are now shown alongside. The three-way row does not reproduce exactly: it is
\textbf{0.9816} here against 0.9957 as first written, and because the original probe's
hyperparameters were never recorded the 0.014 gap cannot be attributed. The reproducible value is
quoted with its method (standardise, logistic regression $C=1.0$, 5-fold stratified out-of-fold
balanced accuracy, 20 size-balanced draws). No $p$ is shown for the three-way rows: ours is floored
at $1/21$ by the number of draws and the original scheme is unrecoverable, so neither is
evidential --- the separation against a null at 0.3345 is. Full account in
\texttt{results/R19-external-site-probe/}.

Two genuinely independent cohorts separate perfectly \emph{under Phikon-v2}.
\textcolor{qred}{\textbf{Qualified 2026-08-17 (R20).}} On the same 76 H\&E slides with the same
probe, \texttt{facebook/dinov2-large} reaches 0.8768 against a null of 0.4990 at $p = 0.0050$. The
\emph{existence} of the signature is encoder-independent and strongly significant either way; its
\emph{totality} is a property of the histology-pretrained encoder, and the word ``perfectly''
belongs to Phikon-v2 rather than to the archive. R20 also records that the 78-slide set used above
was assembled from a partial local staging and omitted 51 of the archive's 80 HTMCP slides ---
Phikon-v2 returns 1.0000 on the full set too, so the conclusion is unaffected, but the provenance
was not documented. We do not present the separation as surprising ---
an Italian population-based series and a sub-Saharan HIV-positive series differ in stain,
scanner and protocol, and a representation with no invariance training will encode that. The
point is the consequence: a model trained on either will have this signal available, and will
use it if the label correlates with it.

\subsection{Federation recovers centralized subtype performance without moving data}

\begin{center}
\begin{tabular}{@{}lcc@{}}
\toprule
\textbf{Regime} & \textbf{Subtype AUROC} & \textbf{95\% CI} \\
\midrule
Local model, evaluated on its own site's patients & 0.9999 & 0.9997--1.0000 \\
Local models, evaluated on the other sites' patients & 0.8372 & 0.8331--0.8412 \\
Federated averaging & \textbf{0.9760} & 0.9674--0.9846 \\
Centralized pooling (upper bound) & 0.9699 & 0.9592--0.9807 \\
\bottomrule
\end{tabular}
\end{center}

A single-site model scores essentially perfectly at home and 0.837 elsewhere. Federated
averaging reaches 0.976, statistically indistinguishable from --- numerically slightly above ---
the centralized 0.970, while transmitting 301\,MiB of weights per full run against
26\,GiB of features that pooling would have moved, an 88-fold reduction.

\textbf{We claim this for subtype only.} On EGFR and on overall survival the centralized-minus-local
gap is itself small, so any ``fraction of gap recovered'' statistic is numerically unstable and
we do not report one. The defensible statement for those endpoints is that federation was never
worse than the best alternative tested, not that it recovered a gap.

\subsection{What a new institution must supply}

On CPTAC (432 patients, an institution set held back until this experiment), eight seeds:

\begin{center}
\begin{tabular}{@{}llrrr@{}}
\toprule
\textbf{Endpoint} & \textbf{Labels used} & \textbf{Mean AUROC} & \textbf{SD} & \textbf{Eval.\ $n$} \\
\midrule
Subtype & 0 (zero-shot) & 0.981 & 0.008 & 173 \\
Subtype & 100 & 0.985 & 0.008 & 173 \\
\midrule
EGFR & 0 (zero-shot) & 0.712 & 0.088 & 36--43 \\
EGFR & 100 & 0.778 & 0.085 & 36--43 \\
\bottomrule
\end{tabular}
\end{center}

For subtype the answer is that a new institution needs \textbf{no labelled cases}: the entire
gain from 0 to 100 labels is 0.004, within a seed standard deviation of 0.008.

For EGFR the honest analysis is a paired one. The seed-to-seed range at fixed label count is
approximately 0.30, seven times the effect of adding labels, because each seed draws its
labelled cases from --- and therefore shrinks --- the evaluation set. Pairing within seed removes
that shared variance: \textbf{7 of 8 seeds improve, mean paired gain $+0.042$}. So more labels
reliably help, while no particular absolute value is well determined. Every EGFR figure here
rests on 36--43 evaluation patients against subtype's 173, and should not be read as comparably
precise. Independently, the development cohort contains 74 EGFR-mutant patients, supporting
detection only of AUROC $\geq 0.584$; weaker effects are invisible by construction.

\subsection{An incidental observation}

Phikon-v2 embeddings of PD-L1 immunohistochemistry slides predict a tumour proportion score of
50\% or above at AUROC 0.870 (95\% CI 0.811--0.922, $n = 163$, 67 positive) with no
immunohistochemistry-specific training. We report this because it was unexpected, not because
this study was designed to establish it.

\section{Discussion}

The practical consequence is a reporting standard. A cross-validated figure from a pooled
multi-institution cohort is an upper bound, and for subtype-like labels whose prevalence varies
between institutions the bound is loose --- here by 17 points. Site-grouped evaluation costs
nothing to implement and should be the default; where it is impossible, the inflation should be
estimated rather than ignored.

Federation is not merely a governance convenience. Because a single-site model scores 0.9999 at
home and 0.837 away, the local figure a hospital would generate for itself is close to
meaningless as a deployment estimate. Federated averaging recovered the pooled result here
while moving 88-fold less data, which makes the argument for consortium training practical
rather than only ethical.

\subsection{Limitations}

\textbf{The federation is simulated.} Weights never crossed a real institutional boundary. The
five silos are \emph{site-anchored partitions of a single collection}, not independent
acquisitions, and they therefore share TCGA's central processing; the 301\,MiB communication
figure is a true measurement of the protocol, not of a deployment.

\textbf{The within-TCGA probe figures are not comparable to each other across groupings.} Five
constructed silos, each pooling several tissue-source sites, give 0.706; the five largest
\emph{individual} tissue-source sites give 0.903. Individual sites are internally more
homogeneous and thus easier to identify. This is a difference of grouping, not of method, and
neither number should be quoted as an improvement on the other.

\textbf{The external cohorts are small and one is single-class.} Seventy-eight H\&E slides from
two cohorts; EAGLE is entirely adenocarcinoma and so cannot test subtype discrimination at all.
They support a qualitative claim about the site signal and nothing quantitative about transfer.

\textbf{EGFR is underpowered throughout} (74 mutants; minimum detectable AUROC 0.584), and every
EGFR statement in this paper is power-qualified.

\textbf{Five encoders now, not one.} \textcolor{qred}{\textbf{Updated 2026-08-18 (R18, R21).}} This
limitation originally read that all conclusions concern Phikon-v2 features. The experiment has since
been run. Re-encoding all 1{,}053 slides with \texttt{facebook/dinov2-large} --- same
\texttt{Dinov2Model} class, same 1024-dimensional width, same 24 layers, same DINOv2
self-supervision, same CLS readout, same ImageNet normalisation, trained on natural images instead
of histology, with byte-identical tile grids --- reproduces the effect and makes it \emph{larger}:
subtype goes 0.7356 site-grouped to 0.9390 random, an inflation of \textbf{+0.2034} against
Phikon-v2's +0.1710. On fraction of remaining headroom the two are 0.769 and 0.852. The site
signature is therefore a property of the archive rather than of histology pretraining. What remains
untested is whether a representation trained with \emph{explicit} stain or site invariance
suppresses it, which is a different and still-open question.

\textbf{But the \emph{magnitude} does not transfer, and that qualifies this paper's headline.}
R21 added three pathology foundation models --- UNI, H-optimus-0 and Virchow2 --- for five in total,
all over byte-identical tiles. Subtype leakage clears the pre-declared bar in every one, so the
phenomenon is universal. Its size is not:

\begin{center}
\small
\begin{tabular}{@{}llrrr@{}}
\toprule
\textbf{Encoder} & \textbf{Corpus} & \textbf{Site-grouped} & \textbf{Inflation} &
\textbf{Relative} \\
\midrule
\texttt{owkin/phikon-v2} & histology & 0.7992 & $+0.1710$ & 0.852 \\
\texttt{facebook/dinov2-large} & natural images & 0.7356 & $+0.2034$ & 0.769 \\
\texttt{MahmoodLab/UNI} & histology & 0.9328 & $+0.0413$ & 0.615 \\
\texttt{bioptimus/H-optimus-0} & histology & 0.9211 & $+0.0473$ & 0.600 \\
\texttt{paige-ai/Virchow2} & histology & \textbf{0.9497} & $+0.0223$ & 0.444 \\
\bottomrule
\end{tabular}
\end{center}

The three newer pathology encoders are far better at the site-disjoint task --- Virchow2 reaches
0.9497 where Phikon-v2 manages 0.7992 --- and show correspondingly smaller inflation. So the
\textbf{0.171} figure this paper leads with should be read as a property of a particular encoder,
not a constant of the archive. The direction of the effect generalises across all five; the
magnitude generalises past none of them. R21 also records, and explicitly declines to claim, an
inverse relationship between site-disjoint capability and relative inflation (Spearman
$\rho = -0.80$, $p = 0.104$, $n = 5$) --- noticed after the fact and reported only as grounds for a
pre-registration.

\textbf{Subtype and KEAP1 only for the fold experiment.} The inflation figure is established for
two labels; its magnitude for others will depend on how those labels distribute across sites.

\section*{Data and code availability}

All imaging is publicly available: TCGA and CPTAC via the NCI Genomic Data Commons,
CDDP\_EAGLE-1 and CGCI-HTMCP-LC likewise. Slide-level manifests with checksums, encoded feature
manifests recording per-slide tile counts and measured microns-per-pixel, analysis
configurations hashed before any label was read, integrity-gate records, and out-of-fold
predictions for every model are retained and available on request.

\section*{Competing interests}

The author is affiliated with a commercial entity developing computational pathology methods.

\end{document}
